\bigskip
\noindent Autres pistes de l'onglet Findings pour ce thème, dans des cotes qui n'ont pas de section ici :\par
\begin{itemize}\small\setlength{\itemsep}{0pt}
\item Cote~138 : critère pour qu'une partie $A$ de la sphère définisse par image réciproque une nouvelle structure de carte : $A$ connexe, et $A$ contient au plus un des points $0$, $1$, $\infty$ ; telle qu'elle est lue, la seconde condition est contredite par l'exemple $A = [0, 1]$, et la lecture de la page doit être vérifiée sur le fac-similé avant toute formalisation. Mathlib n'a pas les cartes sur les surfaces. On se reportera aux notes de Grothendieck (\lot{138}{7}{122}).\piste{138-operations-subset-criterion}
\end{itemize}

\begin{encadre}
\begin{theoreme*}[Grothendieck, pages 122 à 124 et 127, opérations sur les cartes\piste{138-operations-belyi-extending}]\nonverifie
Une opération sur les cartes orientées donnée par une partie $A$ de la sphère de référence, tirée en arrière le long du revêtement $f : X \to S$ de chaque carte, est, sauf un cas exceptionnel, la composition avec un revêtement $\varphi : S \to S$ ramifié seulement au-dessus de $0$, $1$, $\infty$, avec $\varphi^{-1}\{0, 1, \infty\} \supset \{0, 1, \infty\}$ et une ramification $2$ au-dessus de $1$, de sorte que $A = \varphi^{-1}([0, 1])$. La page demande ce qu'il advient dans le cas exceptionnel.
\end{theoreme*}
\end{encadre}
{\small Une carte orientée est un revêtement fini de la sphère non ramifié hors de $0$, $1$, $\infty$, la carte étant l'image réciproque de $[0, 1]$ ; c'est le dictionnaire des cartes et des revêtements de Jones et Singerman (1978) \cite{jonessingerman1978}. Par le théorème de Belyi (1979) \cite{belyi1979}, les courbes qui portent de tels revêtements sont exactement les courbes définies sur $\overline{\mathbf Q}$, et $\mathrm{Gal}(\overline{\mathbf Q}/\mathbf Q)$ opère sur les cartes par son action sur le groupe profini $\hat F_2$ de la section~\ref{sec:145}, qui se factorise par $\widehat{GT}$ (Drinfeld 1990 \cite{drinfeld1990} ; Ihara 1991 \cite{ihara1991}). L'opération de l'encadré est la composition avec un revêtement $\varphi$ du même type, ramifié seulement au-dessus de $0$, $1$, $\infty$. Mathlib n'a pas les revêtements de Belyi de la sphère. On se reportera aux notes de Grothendieck (\lot{138}{7}{122}).\par}
